\section{\ref{sec:eetypes}장. 계측기로서의 뉴런}

개별 세포의 동작 모드는 전극을 통한 전류와 전압 기록으로 직접 측정되었다. 적분기의 수학과 부류 구분은 고전 이론이다. 뇌가 모드 재구성을 계산에 쓴다는 것은 여전히 가설이다.

{\footnotesize
\setlength{\tabcolsep}{3pt}\renewcommand{\arraystretch}{1.15}
\begin{longtable}{>{\raggedright}p{0.5cm}>{\raggedright}p{4.0cm}>{\raggedright}p{5.6cm}>{\raggedright}p{2.3cm}>{\raggedright\arraybackslash}p{1.7cm}}
\toprule
No. & 명제 & 실험과 측정한 것 & 출처 & 상태 \\
\midrule\endfirsthead
\toprule
No. & 명제 & 실험과 측정한 것 & 출처 & 상태 \\
\midrule\endhead
1 & 공진기: 전압 변화에 맞서는 느린 전류가 뉴런을 몇 Hz에서 수십 Hz에 최대가 있는 대역 통과 필터로 만든다(\ref{sec:resonator}절) & 절편의 \nt{GLUT} 뉴런에 주파수가 서서히 오르는 사인파 전류를 넣고 임피던스 $|Z(f)|$를 측정: $33^\circ$C에서 $2$--$5$~Hz에 최대($38^\circ$C에서 약 $7$~Hz). 공진은 느린 칼륨 전류와 양이온 전류가 만들고 지속성 나트륨 전류가 키운다. 리뷰는 같은 기법을 여러 세포 부류에 적용한다 & \cite{HuVervaekeStorm2002,HutcheonYarom2000} & 측정됨 \\
2 & 느린 전류는 인덕턴스처럼 행동하고, 막의 커패시턴스와 함께 공진 회로를 이룬다 & 전류에 대한 축삭의 문턱 아래 반응을 측정하고 선형화한 컨덕턴스 방정식으로 기술했다. 등가 회로에는 값이 전압에 따라 달라지는 “현상론적” 인덕턴스가 들어 있다 & \cite{Mauro1970} & 이론 \\
3 & 되먹임이 있는 무리의 시상수 $\tau_{\text{eff}}=\tau/(1-w)$~\eqref{eq:perfectint}; $\tau=0.1$~s, $w=0.995$에서 $20$~s & 본문에서 무리의 선형 방정식으로 유도. 눈 위치 유지의 메커니즘으로는 이론 연구에서 제안되었다 & 본문에서 유도; \cite{Seung1996} & 이론 \\
4 & 눈 위치 적분기는 개별 세포의 성질이 아니라 네트워크로 입력을 유지한다 & 물고기에서 눈 위치를 유지하는 핵의 세포 안 기록: 빠르게 눈을 돌린 뒤 뉴런의 발화율이 계단처럼 바뀌어 유지된다. 세포 하나에 짧은 전류를 넣으면 잠깐 이동할 뿐이고, 전압 계단은 문턱 아래에서도 유지된다. 이를 유지하는 것은 시냅스 입력이다 & \cite{Aksay2001} & 측정됨 \\
5 & 무누설 적분기에는 학습에 의한 끊임없는 조정이 필요하다. 어긋나면 잊거나 포화로 치닫는다 & 물고기를 눈 위치의 $\pm$에 비례하는 시각 되먹임으로 훈련: 유지가 불안정해지거나 누설되고, 뉴런의 시상수가 한 자릿수 넘게, $<1$~s까지 떨어졌다. 평범한 시각 되먹임은 안정성을 서서히 되돌렸다 & \cite{Major2004} & 측정됨 \\
6 & 쌍안정 뉴런: 짧은 입력이 오래가는 플래토를 켜고 억제가 끈다. 이 성질은 세로토닌이 켠다(\ref{sec:bistable}절) & 고양이에서 근육을 제어하는 \nt{ACET} 뉴런의 세포 안 기록: 짧은 시냅스 흥분이 뉴런을 오래가는 발화로 옮기고 짧은 억제가 재설정한다. 척수를 절단한 고양이에서는 세로토닌 전구체를 주입한 뒤 이 상태가 나타난다 & \cite{Hounsgaard1988} & 측정됨 \\
7 & 두 부류: 적분기(발화율이 0부터 커진다)와 공진기(문턱에서 곧바로 유한한 발화율) & 부류는 분기 이론에서 유도되었다. 실험: 피질 절편에서 \nt{GLUT} 뉴런은 얼마든지 낮은 발화율로 발화를 시작하고(1형), 빠른 \nt{GABA} 뉴런은 유한한 발화율로 도약한다(2형) & \cite{Izhikevich2007,Tateno2004} & 측정됨 \\
8 & 같은 뉴런이 적분기 또는 동시성 검출기가 된다. 억제가 합산 창을 짧게 한다 & 해마 절편에서 흥분성 입력이 문턱까지 더해지는 창은 $2$~ms보다 짧다. 흥분 직후에 오는 피드포워드 억제가 창을 짧게 한다. 억제가 약한 가지돌기에서는 창이 넓다 & \cite{Pouille2001} & 측정됨 \\
9 & 축삭으로 된 지연선(표~\ref{tab:eetypes}) & 원숭이올빼미에서 동시성 검출기로 가는 축삭의 지연을 측정: 축삭 길이가 다르면 지연이 다르고, 검출기는 소리 도달 시간차에 맞춰져 있다 & \cite{Carr1990} & 측정됨 \\
10 & 깨어 있을 때는 페이스메이커, 잘 때는 조용한 세포 & \nt{SERO} 뉴런은 낮에는 거의 규칙적으로 발화하고, 서파 수면에서 느려지며, 렘수면에서는 거의 조용하다(자세한 내용은 \ref{sec:sleep}장) & \cite{JacobsAzmitia1992,McGintyHarper1976} & 측정됨 \\
11 & 소자는 이온 채널과, 그 채널을 조정하는 브로드캐스트 채널이 정한다(명제~\ref{prop:eetypes}) & 작은 네트워크에서 입증: 조절 물질이 뉴런의 고유 성질과 시냅스 세기를 바꿔, 같은 회로가 서로 다른 리듬을 낸다 & \cite{Marder2012} & 측정됨 \\
12 & 뇌는 모드 재구성을 계산에 쓴다(명제~\ref{prop:eetypes}) & 과제를 푸는 데 세포 모드의 재구성이 필요했다는 직접 실험은 없다. 조절 물질이 회로의 출력을 바꾸는 실험만 있다 & \cite{Marder2012} & 가설 \\
\bottomrule
\end{longtable}}
